%
% File ACL2016.tex
%

\documentclass[11pt]{article}
\usepackage{acl2016}
\usepackage{times}
\usepackage{latexsym}
\usepackage{url}
\usepackage{booktabs}
\usepackage{graphicx}
\usepackage{color}
\usepackage{amsmath}
\aclfinalcopy 

\usepackage[authoryear]{natbib}
\usepackage{url}

\title{Project Gutenberg language diversity over time}
\author{Dirk van Wezel \\
S4910583 \\
Length: 3-4 pages}
\date{}

\begin{document}
\maketitle

\begin{abstract}

\end{abstract}


\section{Introduction}
This study analyses the linguistic diversity in Project Gutenberg's digital library collection over time, specifically examining changes in the proportion of English versus non-English books before and after 2010. The research aims to understand how increased internet penetration in non-English speaking countries may have influenced the linguistic composition of digital preservation efforts.
\section{Related Work}
Several studies have examined linguistic diversity in digital libraries and Project Gutenberg specifically:

\cite{leb} explored multilingualism on the web and Project Gutenberg's early efforts at linguistic diversification, providing important historical context for understanding the evolution of digital preservation efforts.

2. Gerlach, M., \& Font-Clos, F. (2020) analysed a curated Project Gutenberg database, revealing a strong bias toward Western languages, with English comprising 81\% of all books. Their work established baseline metrics for understanding linguistic representation in digital libraries.

\section{Data}

\paragraph{Pre-processing}


\section{Predicted Results}

\paragraph{Discussion} 


\section{Conclusion}

\bibliographystyle{chicago}
\bibliography{mybib.bib}

\end{document}



